Soient les formules suivantes~:
\ce{H2O}, \ce{CO}, \ce{O2}, \ce{H2O2}, \ce{H2}, \ce{HCN}, \ce{N2}, \ce{N2O},
\ce{C2H4}, \ce{Cl2}

Classer ces molécules en corps pur simples et corps pur composés~:
\textbf{(10 \texttimes{} 0,25pt)} \par
\begin{minipage}{\linewidth}
\begin{table}[H]
    \centering
    \setlength{\arrayrulewidth}{0.8pt}
    %\setlength{\tabcolsep}{18pt}
    \renewcommand{\arraystretch}{1.5}
    \begin{tabularx}{\textwidth}{|C|C|}
        \hline Corps pur simple & Corps pur composé \\
        \hline \raisebox{1.0cm} & \\ 
        \hline  
    \end{tabularx}
\end{table}
\end{minipage}

